\section{Discussion}\label{sec:discussion}
\paragraph{What CFG changes.} The clearest result is a dose--response one. Raising $w$ from 2 to 15 makes object presence persistent earlier: under a restart that re-noises by eight steps, a trajectory at $w=15$ has settled whether its objects appear several steps before a trajectory at $w=2$ (Section~\ref{sec:resROBUST}). Over the same range, the rate at which the objects appear at all rises and saturates near $w\approx10$. Two common worries do not explain the timing effect. Stronger guidance is known to reduce seed diversity, which would raise agreement between any two images. But chance agreement is computed separately at each $w$, is not higher at $w=15$ on these scenes, and the uncorrected agreement shows the same ordering. The worry that does survive is measurement noise. Correcting $C(t)$ for a judge that errs independently on 14\% of images, the audited rate, removes the difference, so the result is informative only to the extent that the presence judge's errors are content-dependent. Our audit measured how often the judge errs, not how its errors are distributed across seeds. Auditing seed-paired images, which would measure that directly, is the most useful next step for this line of work.

\paragraph{Commitment times are not yet a stable property.} The development split suggested an appealing story: colour assignment is fixed at least as early as object presence and numerosity stays open longest. The available test scenes do not reproduce it. Binding becomes persistent after presence, joint binding about eight steps later than on DEV, and the shift appears under both restart probes. Several factors could produce it, and the partial test cannot separate them: a different scene composition, the DEV precision gate that already failed, the unvalidated binding judge, and the left-censoring of the DEV binding curve at the first grid step. We draw one methodological conclusion: a commitment time estimated on about 100 scenes, with intervals of more than $\pm$7 steps, is too imprecise to schedule an intervention on.

\paragraph{Why CAG does not help, and what the controls show.} CAG was meant to spend guidance while a constraint is still open and to release it afterwards. Two findings explain why it did not work here. First, its switch steps came from DEV commitment times that did not transfer, so on the test scenes it lowers guidance at step 12 while binding is still unsettled. Second, its registered comparison confounds timing with dose. Every CAG arm applies less mean guidance than constant $w=7.5$, and the dose-matched controls attribute the diversity gain entirely to that lower mean: the constant schedule at the same mean is as diverse, and the reversed schedule more so. At equal mean, front-loading guidance is the \emph{least} diverse of the three orders and no better in adherence. This agrees with interval guidance \citep{kynkaanniemi2024interval}, which reports that guidance at the highest noise levels harms diversity, and with step-selection results \citep{grimal2025step}, where strengthening the conditioning late helps alignment more than doing so early. Our data add a controlled, dose-matched comparison on compositional constraints; they do not show that any timed schedule beats constant guidance of the same mean. Online, prompt-dependent schedules \citep{dynamiccfg2025} avoid the transfer problem by measuring during sampling, and they are the natural comparison for a revised CAG.

\paragraph{Attention as a cheap proxy.} The step at which cross-attention maps stop changing orders constraint types almost opposite to persistence (concordance \rbHfConc, $[\rbHfLo,\rbHfHi]$), and within types the scene-level correlations are not distinguishable from zero. On this evidence, attention stabilisation in Stable Diffusion~1.5 is not a substitute for the interventional measurement. This is relevant to methods that read early attention as a forecast of the final image \citep{diffusionprobe2026,guo2024initno}. Their target is image quality or alignment, not the timing of a particular constraint, so our result does not contradict them. It does caution against equating attention that has stopped moving with an outcome that has been decided.

\paragraph{What a reader can take from the null results.} Three negative results are informative because the design anticipated them. The full-resample control and the HA precision gate failed on DEV before any test image existed, which is why HA was demoted to exploratory before it could be over-read. HE has a reversed arm and a matched-mean arm, which is what lets us attribute the diversity effect to dose. And the HD stress test shows how precise a judge must be before timing differences of a few steps become interpretable: with independent errors, even 95\% accuracy leaves the HD interval touching zero.
